% Bölüm 19 — taslak plan. Tam metin yazıldığında bu dosyanın yerini alır.
\chapter{Küme Kuramı ve Bağımsızlık Sonuçlarına Bakış}\label{bol:kume-kurami}

\begin{bolumhedefleri}
\begin{itemize}
  \item ZFC aksiyomlarını
  \item Tip kuramı ile küme kuramı arasındaki ilişkiyi
  \item Süreklilik hipotezinin bağımsızlığının formal ispatını (bakış)
\end{itemize}
\end{bolumhedefleri}

\begin{taslakplan}
\textbf{Planlanan kesimler}
\begin{enumerate}[label=\thechapter.\arabic*, leftmargin=2.4em]
  \item ZFC aksiyomları
  \item Tip kuramında küme modelleri
  \item Sıra ve kardinal sayılar (Mathlib)
  \item Zorlama (forcing) fikri
  \item Flypitch: CH'nin bağımsızlığının formal ispatı
\end{enumerate}

\textbf{Planlanan Lean örnekleri}
\begin{itemize}
  \item \lc{Mathlib'de Cardinal ve Ordinal örnekleri}
\end{itemize}

\textbf{Temel kaynaklar:} \citet{cohen1963,han2020flypitch,mathlib2020}.
\end{taslakplan}

